\section{Cel archiwizacji}\label{sec:cel}
Repozytorium powinno pozwolić innej osobie obejrzeć gotowy PDF i
ponownie złożyć go z tych samych źródeł. Oznacza to zachowanie plików
\texttt{.tex}, bibliografii, danych oraz rysunków, od których dokument
zależy. Plik PDF jest wynikiem, ale nie zastępuje źródła. Z kolei sam
plik \texttt{.tex} bez użytych grafik może przestać się kompilować po
pobraniu projektu.

W opisie projektu należy odróżnić dane przykładowe od pomiarowych,
podać źródła materiałów i wskazać kompilator. Warto używać ścieżek
względnych: \texttt{\textbackslash input\{sections/procedura\}} działa
tak samo w kopii repozytorium jak w folderze autora. W rozdziale
\ref{sec:procedura} pokazano prostą procedurę kontroli.
